\rSec2[expected.ref]{Partial specialization of \tcode{expected} for reference types}

\rSec3[expected.ref.general]{General}

\begin{codeblock}
namespace std {
  template<class T, class E>
  class @\libglobal{expected}@<T&, E> {
  private:
    using @\exposid{error-value-type}@ = remove_cv_t<remove_reference_t<E>>; // \expos

  public:
    using @\libmember{value_type}{expected}@ = T;
    using @\libmember{error_type}{expected}@ = E;
    using @\libmember{unexpected_type}{expected}@ = unexpected<E>;

    template<class U> using rebind = expected<U, error_type>;

    // \ref{expected.ref.cons}, constructors

    expected() = delete;

    constexpr expected(const expected& rhs) noexcept(
        is_nothrow_copy_constructible_v<E>);

    constexpr expected(expected&& rhs) noexcept(@\seebelow@);

    template<class... Args> constexpr expected(in_place_t, Args&&...) = delete;

    template<class U = T>
    constexpr explicit(@\seebelow@)
        expected(U&& u) noexcept(is_nothrow_constructible_v<T&, U>);

    template<class U> constexpr expected(U&&) = delete;

    template<class U, class G>
    constexpr explicit(@\seebelow@) expected(const expected<U&, G>& rhs);

    template<class U, class G>
    constexpr explicit(@\seebelow@) expected(expected<U&, G>&& rhs);

    template<class U, class G>
    constexpr explicit(@\seebelow@) expected(const expected<U&, G>& rhs);

    template<class U, class G>
    constexpr explicit(@\seebelow@) expected(expected<U&, G>&& rhs);

    template<class G> constexpr explicit(@\seebelow@) expected(const unexpected<G>& e);

    template<class G> constexpr explicit(@\seebelow@) expected(unexpected<G>&& e);

    template<class G>
    constexpr explicit(@\seebelow@) expected(const unexpected<G>& e) noexcept;

    template<class G> constexpr explicit(@\seebelow@) expected(unexpected<G>&& e) noexcept;

    template<class G> constexpr expected(const unexpected<G>&) = delete;

    template<class G>
      requires(is_reference_v<E> && !is_reference_v<G>)
    constexpr expected(unexpected<G>&&) = delete;

    template<class... Args> constexpr explicit expected(unexpect_t, Args&&... args);

    template<class... Args> constexpr expected(unexpect_t, Args&&...) = delete;

    template<class... Args>
      requires(is_reference_v<E> && !is_constructible_v<E, Args...> &&
               !@\exposid{unexpect-dangles-v}@<E, Args...>)
    constexpr expected(unexpect_t, Args&&...) = delete;

    template<class U, class... Args>
    constexpr explicit expected(unexpect_t, initializer_list<U> il, Args&&... args);

    template<class U, class... Args>
    constexpr expected(unexpect_t, initializer_list<U>, Args&&...) = delete;

    // \ref{expected.ref.dtor}, destructor

    constexpr ~expected();

    // \ref{expected.ref.assign}, assignment

    constexpr expected& operator=(const expected& rhs) noexcept(
        is_nothrow_copy_constructible_v<E> && is_nothrow_copy_assignable_v<E>);

    constexpr expected& operator=(expected&& rhs) noexcept(@\seebelow@);

    template<class U = T> constexpr expected& operator=(U&& u);

    template<class G> constexpr expected& operator=(const unexpected<G>& e);

    template<class G> constexpr expected& operator=(unexpected<G>&& e);

    template<class G> constexpr expected& operator=(const unexpected<G>& e);

    template<class G> constexpr expected& operator=(unexpected<G>&& e);

    template<class G> constexpr expected& operator=(const unexpected<G>&) = delete;

    template<class G>
      requires(is_reference_v<E> && !is_reference_v<G>)
    constexpr expected& operator=(unexpected<G>&&) = delete;

    template<class U = T> constexpr T& emplace(U&& u) noexcept(@\seebelow@);

    // \ref{expected.ref.swap}, swap

    constexpr void swap(expected& rhs) noexcept(@\seebelow@);

    friend constexpr void swap(expected& x, expected& y) noexcept(noexcept(x.swap(y)));

    // \ref{expected.ref.obs}, observers

    constexpr T* operator->() const noexcept;
    constexpr T& operator*() const noexcept;

    constexpr explicit operator bool() const noexcept;
    constexpr bool has_value() const noexcept;

    constexpr T& value() const&;
    constexpr T& value() &&;

    constexpr const E& error() const& noexcept;
    constexpr E& error() & noexcept;
    constexpr const E&& error() const&& noexcept;
    constexpr E&& error() && noexcept;

    template<class U = remove_cv_t<T>> constexpr remove_cv_t<T> value_or(U&& def) const;

    template<class G = @\exposid{error-value-type}@>
    constexpr @\exposid{error-value-type}@ error_or(G&& def) const&;

    template<class G = @\exposid{error-value-type}@>
    constexpr @\exposid{error-value-type}@ error_or(G&& def) &&;

    // \ref{expected.ref.monadic}, monadic operations

    template<class F> constexpr auto and_then(F&& f) &;
    template<class F> constexpr auto and_then(F&& f) &&;
    template<class F> constexpr auto and_then(F&& f) const&;
    template<class F> constexpr auto and_then(F&& f) const&&;

    template<class F> constexpr auto or_else(F&& f) &;
    template<class F> constexpr auto or_else(F&& f) &&;
    template<class F> constexpr auto or_else(F&& f) const&;
    template<class F> constexpr auto or_else(F&& f) const&&;

    template<class F> constexpr auto transform(F&& f) &;
    template<class F> constexpr auto transform(F&& f) &&;
    template<class F> constexpr auto transform(F&& f) const&;
    template<class F> constexpr auto transform(F&& f) const&&;

    template<class F> constexpr auto transform_error(F&& f) &;
    template<class F> constexpr auto transform_error(F&& f) &&;
    template<class F> constexpr auto transform_error(F&& f) const&;
    template<class F> constexpr auto transform_error(F&& f) const&&;

    // \ref{expected.ref.eq}, equality operators

    template<class T2, class E2>
      requires(!is_void_v<T2>)
    friend constexpr bool operator==(const expected& x, const expected<T2, E2>& y);

    template<class T2>
    friend constexpr bool operator==(const expected& x, const T2& val);

    template<class E2>
    friend constexpr bool operator==(const expected& x, const unexpected<E2>& e);

  private:
    bool @\exposid{has_val}@; // \expos
    union {
      T* @\exposid{val}@;             // \expos
      unexpected<E> @\exposid{unex}@; // \expos
    };
  };
}
\end{codeblock}

\pnum
An object of type \tcode{expected<T\&, E>} either represents a reference to an object of type \tcode{T}, or holds an error. Member \exposid{has_val} indicates whether the object represents a reference. When it represents a reference, member \exposid{val} points to the referenced object, which is not owned by the \tcode{expected} object. Otherwise, the error is \tcode{\exposid{unex}.error()}.

\pnum
The error is held as an \tcode{unexpected<E>}, and not as an \tcode{E}, so that \tcode{E} can be an lvalue reference type: an \tcode{E\&} cannot be a union member, whereas \tcode{unexpected<E\&>} holds a pointer to an external object\iref{expected.un.ref}.

\pnum
A program that instantiates the definition of \tcode{expected<T\&, E>} with a type for the \tcode{E} parameter that is not a valid template argument for \tcode{unexpected} is ill-formed.

\pnum
\tcode{T} shall be an object type that is not an array type, \tcode{in_place_t}, \tcode{unexpect_t}, or a specialization of \tcode{unexpected}. If \tcode{E} is not an lvalue reference type, \tcode{E} shall meet the \oldconcept{Destructible} requirements (\tref{cpp17.destructible}).


\rSec3[expected.ref.cons]{Constructors}

\indexlibraryctor{expected}%
\begin{itemdecl}
expected() = delete;
\end{itemdecl}

\begin{itemdescr}
\pnum
\remarks
\tcode{expected<T\&, E>} has no default constructor: a reference cannot be null, so there is no empty state to default-construct into.
\end{itemdescr}

\indexlibraryctor{expected}%
\begin{itemdecl}
constexpr expected(const expected& rhs) noexcept(is_nothrow_copy_constructible_v<E>);
constexpr expected(expected&& rhs) noexcept(@\seebelow@);
\end{itemdecl}

\begin{itemdescr}

\pnum
\effects
If \tcode{rhs.has_value()} is \tcode{true}, initializes \exposid{val} with \tcode{rhs.\exposid{val}}, so that \tcode{*this} and \tcode{rhs} refer to the same object; otherwise, initializes \exposid{unex} with \tcode{rhs.\exposid{unex}}.

\pnum
\ensures
\tcode{rhs.has_value() == this->has_value()}.

\pnum
\remarks
The exception specification of the move constructor is equivalent to \tcode{is_nothrow_move_constructible_v<E>}. Each of these constructors is trivial if the corresponding constructor of \tcode{E} is trivial. The copy constructor is defined as deleted unless \tcode{is_copy_constructible_v<E>} is \tcode{true}, and the move constructor unless \tcode{is_move_constructible_v<E>} is \tcode{true}.
\end{itemdescr}

\indexlibraryctor{expected}%
\begin{itemdecl}
template<class... Args>
  constexpr expected(in_place_t, Args&&...) = delete;
\end{itemdecl}

\begin{itemdescr}
\pnum
\remarks
\tcode{expected<T\&, E>} has no in-place value constructor: \tcode{T\&} cannot be constructed in-place. Pass a \tcode{U} convertible to \tcode{T\&} instead.
\end{itemdescr}

\indexlibraryctor{expected}%
\begin{itemdecl}
template<class U = T>
  constexpr explicit(@\seebelow@) expected(U&& u) noexcept(is_nothrow_constructible_v<T&, U>);
\end{itemdecl}

\begin{itemdescr}
\pnum
\constraints
\tcode{remove_cvref_t<U>} is not \tcode{in_place_t}, \tcode{expected}, or a specialization of \tcode{unexpected}; \tcode{is_constructible_v<T\&, U>} is \tcode{true}; and \tcode{reference_constructs_from_temporary_v<T\&, U>} is \tcode{false}.

\pnum
\effects
Let \tcode{r} be the lvalue result of \tcode{T\& r = std::forward<U>(u);}. Initializes \exposid{val} with \tcode{addressof(r)}.

\pnum
\ensures
\tcode{has_value()} is \tcode{true}.
\end{itemdescr}

\indexlibraryctor{expected}%
\begin{itemdecl}
template<class U>
  constexpr expected(U&&) = delete;
\end{itemdecl}

\begin{itemdescr}
\pnum
\constraints
\tcode{reference_constructs_from_temporary_v<T\&, U>} is \tcode{true}.

\pnum
\remarks
A constructor for which \tcode{reference_constructs_from_temporary_v<T\&, U>} is \tcode{true} — one that would bind \tcode{T\&} to a temporary — is defined as deleted.
\end{itemdescr}

\indexlibraryctor{expected}%
\begin{itemdecl}
template<class U, class G>
  constexpr explicit(@\seebelow@) expected(const expected<U&, G>& rhs);
template<class U, class G>
  constexpr explicit(@\seebelow@) expected(expected<U&, G>&& rhs);
\end{itemdecl}

\begin{itemdescr}
\pnum
\constraints
\tcode{is_constructible_v<T\&, U\&>} is \tcode{true}; \tcode{reference_constructs_from_temporary_v<T\&, U\&>} is \tcode{false}; and \tcode{is_constructible_v<E, const G\&>} is \tcode{true}.

\pnum
\effects
If \tcode{rhs.has_value()} is \tcode{true}, initializes \exposid{val} with \tcode{addressof(*rhs)}, so that \tcode{*this} refers to the object referred to by \tcode{rhs}; otherwise, initializes \exposid{unex} with the error of \tcode{rhs}. No object referred to by \tcode{rhs} is moved from.
\end{itemdescr}

\indexlibraryctor{expected}%
\begin{itemdecl}
template<class U, class G>
  constexpr explicit(@\seebelow@) expected(const expected<U&, G>& rhs);
template<class U, class G>
  constexpr explicit(@\seebelow@) expected(expected<U&, G>&& rhs);
\end{itemdecl}

\begin{itemdescr}
\pnum
\constraints
\tcode{is_constructible_v<T\&, U\&>} is \tcode{true}; \tcode{reference_constructs_from_temporary_v<T\&, U\&>} is \tcode{false}; \tcode{G} is a reference type; and \tcode{is_convertible_v<G, E>} is \tcode{true}.

\pnum
\effects
If \tcode{rhs.has_value()} is \tcode{true}, initializes \exposid{val} with \tcode{addressof(*rhs)}, so that \tcode{*this} refers to the object referred to by \tcode{rhs}; otherwise, initializes \exposid{unex} with the error of \tcode{rhs}. No object referred to by \tcode{rhs} is moved from.

\pnum
\remarks
Unlike the value-\tcode{E} overload above, this overload participates in overload resolution only when \tcode{E} and \tcode{G} are both reference types, so the referenced error object is never copied.
\end{itemdescr}

\indexlibraryctor{expected}%
\begin{itemdecl}
template<class G>
  constexpr explicit(@\seebelow@) expected(const unexpected<G>& e);
template<class G>
  constexpr explicit(@\seebelow@) expected(unexpected<G>&& e);
\end{itemdecl}

\begin{itemdescr}
\pnum
\constraints
\tcode{is_constructible_v<E, const G\&>} is \tcode{true}.

\pnum
\effects
Initializes \exposid{unex} with the error of \tcode{e}.

\pnum
\ensures
\tcode{has_value()} is \tcode{false}.
\end{itemdescr}

\indexlibraryctor{expected}%
\begin{itemdecl}
template<class G>
  constexpr explicit(@\seebelow@) expected(const unexpected<G>& e) noexcept;
template<class G>
  constexpr explicit(@\seebelow@) expected(unexpected<G>&& e) noexcept;
\end{itemdecl}

\begin{itemdescr}
\pnum
\constraints
\tcode{is_reference_v<G>} is \tcode{true}; \tcode{is_constructible_v<E, G>} is \tcode{true}; and \tcode{reference_constructs_from_temporary_v<E, G>} is \tcode{false}.

\pnum
\effects
Initializes \exposid{unex} with the error of \tcode{e}.

\pnum
\ensures
\tcode{has_value()} is \tcode{false}.

\pnum
\remarks
This constructor never throws: the referent is bound, not copied.
\end{itemdescr}

\indexlibraryctor{expected}%
\begin{itemdecl}
template<class G>
  constexpr expected(const unexpected<G>&) = delete;
\end{itemdecl}

\begin{itemdescr}
\pnum
\constraints
\tcode{is_reference_v<E>} is \tcode{true} and \tcode{is_reference_v<G>} is \tcode{false}.

\pnum
\remarks
When \tcode{E} is a reference type, an overload taking \tcode{unexpected<G>} for a non-reference \tcode{G} is defined as deleted: the referent would live inside the (possibly temporary) source \tcode{unexpected<G>} object, and binding \tcode{E\&} to it would dangle. Use an \tcode{unexpected<E\&>} holding an external object instead.
\end{itemdescr}

\indexlibraryctor{expected}%
\begin{itemdecl}
template<class... Args>
  constexpr explicit expected(unexpect_t, Args&&... args);
\end{itemdecl}

\begin{itemdescr}
\pnum
\constraints
\tcode{is_constructible_v<E, Args...>} is \tcode{true}.

\pnum
\effects
Direct-non-list-initializes \exposid{unex} with \tcode{in_place} and \tcode{std::forward<Args>(args)...}.

\pnum
\ensures
\tcode{has_value()} is \tcode{false}.
\end{itemdescr}

\indexlibraryctor{expected}%
\begin{itemdecl}
template<class... Args>
  constexpr expected(unexpect_t, Args&&...) = delete;
\end{itemdecl}

\begin{itemdescr}
\pnum
\constraints
\tcode{\exposid{unexpect-dangles-v}<E, Args...>} is \tcode{true}.

\pnum
\remarks
When \tcode{E} is a reference type, an overload with the same parameter types is defined as deleted if the single argument would bind \tcode{E\&} to a temporary, or if it is otherwise not usable to construct \tcode{E}.
\end{itemdescr}

\indexlibraryctor{expected}%
\begin{itemdecl}
template<class U, class... Args>
  constexpr explicit expected(unexpect_t, initializer_list<U> il, Args&&... args);
\end{itemdecl}

\begin{itemdescr}
\pnum
\constraints
\tcode{is_reference_v<E>} is \tcode{false}, and \tcode{is_constructible_v<E, initializer_list<U>\&, Args...>} is \tcode{true}.

\pnum
\effects
Direct-non-list-initializes \exposid{unex} with \tcode{in_place}, \tcode{il}, and \tcode{std::forward<Args>(args)...}.

\pnum
\ensures
\tcode{has_value()} is \tcode{false}.

\pnum
\remarks
An overload with the same parameter types is defined as deleted when \tcode{E} is an lvalue reference type. An initializer list cannot provide the required long-lived error referent.
\end{itemdescr}

\indexlibraryctor{expected}%
\begin{itemdecl}
template<class U, class... Args>
  constexpr expected(unexpect_t, initializer_list<U>, Args&&...) = delete;
\end{itemdecl}

\begin{itemdescr}
\pnum
\constraints
\tcode{is_reference_v<E>} is \tcode{true}.

\pnum
\remarks
An overload with the same parameter types is defined as deleted when \tcode{E} is an lvalue reference type. An initializer list cannot provide the required long-lived error referent.
\end{itemdescr}

\rSec3[expected.ref.dtor]{Destructor}

\indexlibrarydtor{expected}%
\begin{itemdecl}
constexpr ~expected();
\end{itemdecl}

\begin{itemdescr}
\pnum
\effects
If \tcode{has_value()} is \tcode{false}, destroys \exposid{unex}. \tcode{T} is not destroyed; \tcode{*this} never owns the object it refers to.

\pnum
\remarks
This destructor is trivial if \tcode{E} is trivially destructible.
\end{itemdescr}

\rSec3[expected.ref.assign]{Assignment}

\indexlibrarymember{operator=}{expected}%
\begin{itemdecl}
constexpr expected& operator=(const expected& rhs) noexcept(is_nothrow_copy_constructible_v<E> &&
                                                            is_nothrow_copy_assignable_v<E>);
constexpr expected& operator=(expected&& rhs) noexcept(@\seebelow@);
\end{itemdecl}

\begin{itemdescr}

\pnum
\effects
If \tcode{rhs.has_value()} is \tcode{true}: if \tcode{has_value()} is \tcode{true}, assigns \tcode{rhs.\exposid{val}} to \exposid{val}; otherwise destroys \exposid{unex} and initializes \exposid{val} with \tcode{rhs.\exposid{val}}. If \tcode{rhs.has_value()} is \tcode{false}, the error of \tcode{rhs} is assigned to or used to initialize \exposid{unex}, as for the primary template. In every case \tcode{*this} comes to refer to the object \tcode{rhs} refers to, or to hold the error of \tcode{rhs}.

\pnum
\returns
\tcode{*this}.

\pnum
\remarks
The exception specification of the move assignment operator is equivalent to \tcode{is_nothrow_move_constructible_v<E> \&\& is_nothrow_move_assignable_v<E>}. The copy assignment operator is defined as deleted unless \tcode{is_copy_assignable_v<E>} is \tcode{true} and \tcode{is_copy_constructible_v<E>} is \tcode{true}, and the move assignment operator unless \tcode{is_move_assignable_v<E>} is \tcode{true} and \tcode{is_move_constructible_v<E>} is \tcode{true}. Each of these operators is trivial if the corresponding constructor and assignment operator of \tcode{E} are trivial and \tcode{is_trivially_destructible_v<E>} is \tcode{true}.
\end{itemdescr}

\indexlibrarymember{operator=}{expected}%
\begin{itemdecl}
template<class U = T>
  constexpr expected& operator=(U&& u);
\end{itemdecl}

\begin{itemdescr}
\pnum
\constraints
\tcode{remove_cvref_t<U>} is neither \tcode{expected} nor a specialization of \tcode{unexpected}, \tcode{is_constructible_v<T\&, U>} is \tcode{true}, and \tcode{reference_constructs_from_temporary_v<T\&, U>} is \tcode{false}.

\pnum
\effects
Let \tcode{r} be the lvalue result of \tcode{T\& r = std::forward<U>(u);}. If \tcode{has_value()} is \tcode{true}, assigns \tcode{addressof(r)} to \exposid{val}. Otherwise, destroys \exposid{unex}, initializes \exposid{val} with \tcode{addressof(r)}, and sets \exposid{has_val} to \tcode{true}; if binding \tcode{r} throws, \tcode{*this} is left unchanged.

\pnum
\returns
\tcode{*this}.
\end{itemdescr}

\indexlibrarymember{operator=}{expected}%
\begin{itemdecl}
template<class G>
  constexpr expected& operator=(const unexpected<G>& e);
template<class G>
  constexpr expected& operator=(unexpected<G>&& e);
\end{itemdecl}

\begin{itemdescr}
\pnum
\constraints
\tcode{is_constructible_v<E, const G\&>} is \tcode{true} and \tcode{is_assignable_v<E\&, const G\&>} is \tcode{true}.

\pnum
\effects
Makes \tcode{*this} hold the error of \tcode{e}, reinitializing \exposid{unex} from \tcode{e} rather than assigning through it.

\pnum
\returns
\tcode{*this}.
\end{itemdescr}

\indexlibrarymember{operator=}{expected}%
\begin{itemdecl}
template<class G>
  constexpr expected& operator=(const unexpected<G>& e);
template<class G>
  constexpr expected& operator=(unexpected<G>&& e);
\end{itemdecl}

\begin{itemdescr}
\pnum
\constraints
\tcode{is_reference_v<G>} is \tcode{true}; \tcode{is_constructible_v<E, G>} is \tcode{true}; and \tcode{reference_constructs_from_temporary_v<E, G>} is \tcode{false}.

\pnum
\effects
Makes \tcode{*this} hold the error of \tcode{e}, reinitializing \exposid{unex} from \tcode{e} rather than assigning through it. \tcode{\exposid{unex}.error()} thereafter refers to the same object as \tcode{e.error()}; the previously referenced object, if any, is not modified.

\pnum
\returns
\tcode{*this}.

\pnum
\remarks
This operator never throws: the referent is bound, not copied.
\end{itemdescr}

\indexlibrarymember{operator=}{expected}%
\begin{itemdecl}
template<class G>
  constexpr expected& operator=(const unexpected<G>&) = delete;
\end{itemdecl}

\begin{itemdescr}
\pnum
\constraints
\tcode{is_reference_v<E>} is \tcode{true} and \tcode{is_reference_v<G>} is \tcode{false}.

\pnum
\remarks
When \tcode{E} is a reference type, an overload taking \tcode{unexpected<G>} for a non-reference \tcode{G} is defined as deleted: it would rebind \exposid{unex} to \tcode{unexpected<G>}'s temporary storage.
\end{itemdescr}

\indexlibrarymember{emplace}{expected}%
\begin{itemdecl}
template<class U = T>
  constexpr T& emplace(U&& u) noexcept(@\seebelow@);
\end{itemdecl}

\begin{itemdescr}
\pnum
\constraints
\tcode{is_constructible_v<T\&, U>} is \tcode{true} and \tcode{reference_constructs_from_temporary_v<T\&, U>} is \tcode{false}.

\pnum
\effects
Rebinds \tcode{*this} to refer to the object bound by \tcode{T\& r = std::forward<U>(u);}: if \tcode{has_value()} is \tcode{false}, destroys \exposid{unex} first. Sets \exposid{val} to \tcode{addressof(r)} and \exposid{has_val} to \tcode{true}.

\pnum
\returns
\tcode{*\exposid{val}}.
\end{itemdescr}

\pnum
Assignment rebinds: assigning to an \tcode{expected<T\&, E>} that holds a value changes which object it refers to. It never assigns through to the referent.

\rSec3[expected.ref.swap]{Swap}

\indexlibrarymember{swap}{expected}%
\begin{itemdecl}
constexpr void swap(expected& rhs) noexcept(@\seebelow@);
\end{itemdecl}

\begin{itemdescr}
\pnum
\constraints
\tcode{is_reference_v<E> || is_swappable_v<E>} is \tcode{true} and \tcode{is_move_constructible_v<E>} is \tcode{true}.

\pnum
\effects
Exchanges the states of \tcode{*this} and \tcode{rhs}. When both hold values, exchanges \exposid{val} and \tcode{rhs.\exposid{val}} — the referenced objects are not swapped. Otherwise behaves as the primary template's \tcode{swap} does for the error.

\pnum
\remarks
The exception specification is equivalent to \tcode{is_nothrow_move_constructible_v<E> \&\& (is_reference_v<E> || is_nothrow_swappable_v<E>)}.
\end{itemdescr}

\indexlibraryglobal{swap}%
\begin{itemdecl}
friend constexpr void swap(expected& x, expected& y) noexcept(noexcept(x.swap(y)));
\end{itemdecl}

\begin{itemdescr}
\pnum
\constraints
\tcode{is_reference_v<E> || is_swappable_v<E>} is \tcode{true} and \tcode{is_move_constructible_v<E>} is \tcode{true}.

\pnum
\effects
Equivalent to \tcode{x.swap(y)}.
\end{itemdescr}

\rSec3[expected.ref.obs]{Observers}

\indexlibrarymember{operator->}{expected}%
\begin{itemdecl}
constexpr T* operator->() const noexcept;
\end{itemdecl}

\begin{itemdescr}
\pnum
\expects
\tcode{has_value()} is \tcode{true}.

\pnum
\returns
\exposid{val}.

\pnum
\remarks
This is a \tcode{const} member function that returns a non-\tcode{const} \tcode{T*}; the constness of \tcode{*this} does not propagate to the referenced object. For deep \tcode{const}, use \tcode{expected<const T\&, E>}.
\end{itemdescr}

\indexlibrarymember{operator*}{expected}%
\begin{itemdecl}
constexpr T& operator*() const noexcept;
\end{itemdecl}

\begin{itemdescr}
\pnum
\expects
\tcode{has_value()} is \tcode{true}.

\pnum
\returns
\tcode{*\exposid{val}}.

\pnum
\remarks
This is a \tcode{const} member function that returns a non-\tcode{const} \tcode{T\&}; the constness of \tcode{*this} does not propagate to the referenced object. For deep \tcode{const}, use \tcode{expected<const T\&, E>}.
\end{itemdescr}

\indexlibrarymember{operator bool}{expected}%
\indexlibrarymember{has_value}{expected}%
\begin{itemdecl}
constexpr explicit operator bool() const noexcept;
constexpr bool has_value() const noexcept;
\end{itemdecl}

\begin{itemdescr}
\pnum
\returns
\exposid{has_val}.
\end{itemdescr}

\indexlibrarymember{value}{expected}%
\begin{itemdecl}
constexpr T& value() const&;
\end{itemdecl}

\begin{itemdescr}
\pnum
\mandates
\tcode{is_copy_constructible_v<\exposid{error-value-type}>} is \tcode{true}.

\pnum
\returns
\tcode{*\exposid{val}} if \tcode{has_value()} is \tcode{true}.

\pnum
\throws
\tcode{bad_expected_access(as_const(error()))} if \tcode{has_value()} is \tcode{false}.
\end{itemdescr}

\indexlibrarymember{value}{expected}%
\begin{itemdecl}
constexpr T& value() &&;
\end{itemdecl}

\begin{itemdescr}
\pnum
\returns
\tcode{*\exposid{val}} if \tcode{has_value()} is \tcode{true}.

\pnum
\throws
\tcode{bad_expected_access(std::move(error()))} if \tcode{has_value()} is \tcode{false}.
\end{itemdescr}

\indexlibrarymember{error}{expected}%
\begin{itemdecl}
constexpr const E& error() const& noexcept;
constexpr E& error() & noexcept;
\end{itemdecl}

\begin{itemdescr}
\pnum
\expects
\tcode{has_value()} is \tcode{false}.

\pnum
\returns
\tcode{\exposid{unex}.error()}.
\end{itemdescr}

\indexlibrarymember{error}{expected}%
\begin{itemdecl}
constexpr const E&& error() const&& noexcept;
constexpr E&& error() && noexcept;
\end{itemdecl}

\begin{itemdescr}
\pnum
\expects
\tcode{has_value()} is \tcode{false}.

\pnum
\returns
\tcode{std::move(\exposid{unex}).error()}.
\end{itemdescr}

\indexlibrarymember{value_or}{expected}%
\begin{itemdecl}
template<class U = remove_cv_t<T>>
  constexpr remove_cv_t<T> value_or(U&& def) const;
\end{itemdecl}

\begin{itemdescr}
\pnum
\constraints
\tcode{is_object_v<T>} is \tcode{true} and \tcode{is_array_v<T>} is \tcode{false}.

\pnum
\mandates
\tcode{is_convertible_v<T\&, remove_cv_t<T>>} and \tcode{is_convertible_v<U, remove_cv_t<T>>} are \tcode{true}.

\pnum
\returns
\tcode{has_value() ? static_cast<remove_cv_t<T>>(*\exposid{val}) : static_cast<remove_cv_t<T>>(std::forward<U>(def))}. The result is an object, never a reference.
\end{itemdescr}

\indexlibrarymember{error_or}{expected}%
\begin{itemdecl}
template<class G = @\exposid{error-value-type}@>
  constexpr @\exposid{error-value-type}@ error_or(G&& def) const&;
\end{itemdecl}

\begin{itemdescr}
\pnum
\constraints
\tcode{is_copy_constructible_v<remove_cv_t<remove_reference_t<E>>>} is \tcode{true} and \tcode{is_convertible_v<G, remove_cv_t<remove_reference_t<E>>>} is \tcode{true}.

\pnum
\returns
\tcode{std::forward<G>(def)} if \tcode{has_value()} is \tcode{true}, \tcode{error()} otherwise. The result is an object, never a reference.
\end{itemdescr}

\indexlibrarymember{error_or}{expected}%
\begin{itemdecl}
template<class G = @\exposid{error-value-type}@>
  constexpr @\exposid{error-value-type}@ error_or(G&& def) &&;
\end{itemdecl}

\begin{itemdescr}
\pnum
\constraints
\tcode{is_move_constructible_v<remove_cv_t<remove_reference_t<E>>>} is \tcode{true} and \tcode{is_convertible_v<G, remove_cv_t<remove_reference_t<E>>>} is \tcode{true}.

\pnum
\returns
\tcode{std::forward<G>(def)} if \tcode{has_value()} is \tcode{true}, \tcode{std::move(error())} otherwise. The result is an object, never a reference.
\end{itemdescr}

\rSec3[expected.ref.monadic]{Monadic operations}

\indexlibrarymember{and_then}{expected}%
\begin{itemdecl}
template<class F>
  constexpr auto and_then(F&& f) &;
template<class F>
  constexpr auto and_then(F&& f) const&;
\end{itemdecl}

\begin{itemdescr}
\pnum
\constraints
\tcode{is_constructible_v<E, decltype(error())>} is \tcode{true}.

\pnum
\mandates
\tcode{remove_cvref_t<invoke_result_t<F, T\&>>} is a specialization of \tcode{expected} and its \tcode{error_type} is the same type as \tcode{E}.

\pnum
\effects
Equivalent to: \tcode{if (has_value()) return invoke(std::forward<F>(f), *\exposid{val}); else return U(unexpect, error());} where \tcode{U} is \tcode{remove_cvref_t<invoke_result_t<F, T\&>>}.

\pnum
\remarks
The member templates \tcode{and_then}, \tcode{or_else}, \tcode{transform}, and \tcode{transform_error} behave as specified for the primary template, with one difference: the value is passed to the callable as \tcode{T\&} for every ref-qualification of \tcode{*this}. An rvalue \tcode{expected<T\&, E>} does not pass its referent as an rvalue; the object referred to is never moved from by these operations.
\end{itemdescr}

\indexlibrarymember{and_then}{expected}%
\begin{itemdecl}
template<class F>
  constexpr auto and_then(F&& f) &&;
template<class F>
  constexpr auto and_then(F&& f) const&&;
\end{itemdecl}

\begin{itemdescr}
\pnum
\constraints
\tcode{is_constructible_v<E, decltype(std::move(error()))>} is \tcode{true}.

\pnum
\mandates
\tcode{remove_cvref_t<invoke_result_t<F, T\&>>} is a specialization of \tcode{expected} and its \tcode{error_type} is the same type as \tcode{E}.

\pnum
\effects
Equivalent to: \tcode{if (has_value()) return invoke(std::forward<F>(f), *\exposid{val}); else return U(unexpect, std::move(error()));} where \tcode{U} is \tcode{remove_cvref_t<invoke_result_t<F, T\&>>}.
\end{itemdescr}

\indexlibrarymember{or_else}{expected}%
\begin{itemdecl}
template<class F>
  constexpr auto or_else(F&& f) &;
template<class F>
  constexpr auto or_else(F&& f) const&;
\end{itemdecl}

\begin{itemdescr}
\pnum
\mandates
\tcode{remove_cvref_t<invoke_result_t<F, decltype(error())>>} is \tcode{expected<T\&, X>} for some type \tcode{X}.

\pnum
\effects
Equivalent to: \tcode{if (has_value()) return G(*\exposid{val}); else return invoke(std::forward<F>(f), error());} where \tcode{G} is \tcode{remove_cvref_t<invoke_result_t<F, decltype(error())>>}.
\end{itemdescr}

\indexlibrarymember{or_else}{expected}%
\begin{itemdecl}
template<class F>
  constexpr auto or_else(F&& f) &&;
template<class F>
  constexpr auto or_else(F&& f) const&&;
\end{itemdecl}

\begin{itemdescr}
\pnum
\mandates
\tcode{remove_cvref_t<invoke_result_t<F, decltype(std::move(error()))>>} is \tcode{expected<T\&, X>} for some type \tcode{X}.

\pnum
\effects
Equivalent to: \tcode{if (has_value()) return G(*\exposid{val}); else return invoke(std::forward<F>(f), std::move(error()));} where \tcode{G} is \tcode{remove_cvref_t<invoke_result_t<F, decltype(std::move(error()))>>}.
\end{itemdescr}

\indexlibrarymember{transform}{expected}%
\begin{itemdecl}
template<class F>
  constexpr auto transform(F&& f) &;
template<class F>
  constexpr auto transform(F&& f) const&;
\end{itemdecl}

\begin{itemdescr}
\pnum
\constraints
\tcode{is_constructible_v<E, decltype(error())>} is \tcode{true}.

\pnum
\effects
Equivalent to: \tcode{if (!has_value()) return U(unexpect, error()); else return expected<U2, E>(in_place, invoke(std::forward<F>(f), *\exposid{val}));} where \tcode{U2} is \tcode{remove_cv_t<invoke_result_t<F, T\&>>} and \tcode{U} is \tcode{expected<U2, E>}.
\end{itemdescr}

\indexlibrarymember{transform}{expected}%
\begin{itemdecl}
template<class F>
  constexpr auto transform(F&& f) &&;
template<class F>
  constexpr auto transform(F&& f) const&&;
\end{itemdecl}

\begin{itemdescr}
\pnum
\constraints
\tcode{is_constructible_v<E, decltype(std::move(error()))>} is \tcode{true}.

\pnum
\effects
Equivalent to: \tcode{if (!has_value()) return U(unexpect, std::move(error())); else return expected<U2, E>(in_place, invoke(std::forward<F>(f), *\exposid{val}));} where \tcode{U2} is \tcode{remove_cv_t<invoke_result_t<F, T\&>>} and \tcode{U} is \tcode{expected<U2, E>}.
\end{itemdescr}

\indexlibrarymember{transform_error}{expected}%
\begin{itemdecl}
template<class F>
  constexpr auto transform_error(F&& f) &;
template<class F>
  constexpr auto transform_error(F&& f) const&;
\end{itemdecl}

\begin{itemdescr}
\pnum
\mandates
\begin{itemize}
\item \tcode{is_object_v<G>} is \tcode{true},
\item \tcode{is_array_v<G>} is \tcode{false},
\item \tcode{is_same_v<G, remove_cv_t<G>>} is \tcode{true},
\item \tcode{\exposid{is-unexpected-specialization}<G>::value} is \tcode{false}.
\end{itemize}

\pnum
\effects
Equivalent to: \tcode{if (has_value()) return G(*\exposid{val}); else return expected<T\&, G2>(unexpect, invoke(std::forward<F>(f), error()));} where \tcode{G2} is \tcode{remove_cv_t<invoke_result_t<F, decltype(error())>>} and \tcode{G} is \tcode{expected<T\&, G2>}.
\end{itemdescr}

\indexlibrarymember{transform_error}{expected}%
\begin{itemdecl}
template<class F>
  constexpr auto transform_error(F&& f) &&;
template<class F>
  constexpr auto transform_error(F&& f) const&&;
\end{itemdecl}

\begin{itemdescr}
\pnum
\mandates
\begin{itemize}
\item \tcode{is_object_v<G>} is \tcode{true},
\item \tcode{is_array_v<G>} is \tcode{false},
\item \tcode{is_same_v<G, remove_cv_t<G>>} is \tcode{true},
\item \tcode{\exposid{is-unexpected-specialization}<G>::value} is \tcode{false}.
\end{itemize}

\pnum
\effects
Equivalent to: \tcode{if (has_value()) return G(*\exposid{val}); else return expected<T\&, G2>(unexpect, invoke(std::forward<F>(f), std::move(error())));} where \tcode{G2} is \tcode{remove_cv_t<invoke_result_t<F, decltype(std::move(error()))>>} and \tcode{G} is \tcode{expected<T\&, G2>}.
\end{itemdescr}

\rSec3[expected.ref.eq]{Equality operators}

\indexlibraryglobal{operator==}%
\begin{itemdecl}
template<class T2, class E2>
  requires(!is_void_v<T2>)
  friend constexpr bool operator==(const expected& x, const expected<T2, E2>& y);
\end{itemdecl}

\begin{itemdescr}
\pnum
\mandates
The expression \tcode{*x == *y} is well-formed and its result is convertible to \tcode{bool}. The expression \tcode{x.error() == y.error()} is well-formed and its result is convertible to \tcode{bool}.

\pnum
\returns
If \tcode{x.has_value() != y.has_value()}, \tcode{false}; otherwise, if \tcode{x.has_value()} is \tcode{true}, \tcode{*x == *y}; otherwise \tcode{x.error() == y.error()}.

\pnum
\remarks
The equality operators behave as specified for the primary template, comparing referents through \tcode{operator*} and errors through \tcode{error()}.
\end{itemdescr}

\indexlibraryglobal{operator==}%
\begin{itemdecl}
template<class T2>
  friend constexpr bool operator==(const expected& x, const T2& val);
\end{itemdecl}

\begin{itemdescr}
\pnum
\constraints
\tcode{\exposid{is-expected-specialization}<T2>::value} is \tcode{false}.

\pnum
\mandates
\tcode{T2} is not a specialization of \tcode{expected}. The expression \tcode{*x == val} is well-formed and its result is convertible to \tcode{bool}.

\pnum
\returns
\tcode{x.has_value() \&\& static_cast<bool>(*x == val)}.
\end{itemdescr}

\indexlibraryglobal{operator==}%
\begin{itemdecl}
template<class E2>
  friend constexpr bool operator==(const expected& x, const unexpected<E2>& e);
\end{itemdecl}

\begin{itemdescr}
\pnum
\mandates
The expression \tcode{x.error() == e.error()} is well-formed and its result is convertible to \tcode{bool}.

\pnum
\returns
\tcode{!x.has_value() \&\& static_cast<bool>(x.error() == e.error())}.
\end{itemdescr}
